% P11 paper module: R27 constrained Gamma Mosco limit and fixed-window resolvent reduction.

\subsection{Constrained Gamma Mosco limit and fixed-window resolvent reduction}
\label{subsec:o3z-constrained-gamma}

The odd terminal metrics diverge on every fixed nonzero smooth odd source vector, but
R17 exhibits moving near-null vectors converging to higher-jet directions while their
future Schur energy vanishes.  These two facts fit into a single variational limit: the
first boundary functional becomes a hard constraint and the surviving finite energy is
the Gamma form.

Fix a source level $X>0$ and write
\[
 \mathcal H_X:=\KXR{X}^{-},
 \qquad
 a_{X,U}(f):=q_U^X(J_{X,U}f)
 \quad(U>X).
\]
Let
\[
 H_X^0:=\ker\beta_X^{(0)}\cap\mathcal H_X.
\]
Since $\beta_X^{(0)}$ is an $L^2$ pairing with a fixed bounded kernel and the graph norm
dominates $L^2$, $H_X^0$ is a closed hyperplane.

\begin{theorem}[Constrained Gamma Mosco limit]
\label{thm:o3z-mosco}
As $U\to\infty$, the finite terminal forms $a_{X,U}$ Mosco-converge on
$\mathcal H_X$ to the closed extended quadratic form
\[
 \boxed{
 a_{X,\infty}(f)
 :=
 \begin{cases}
 \mathfrak c_{\Gamma,X}[f],&f\in H_X^0,\\[2pt]
 +\infty,&f\notin H_X^0.
 \end{cases}
 }
\tag{O3Z.1}\label{eq:o3z-limit-form}
\]
\end{theorem}

\begin{proof}
We first prove the weak liminf.  On the old source window,
\eqref{eq:PhiT-asymptotic} implies the uniform $L^2(-X,X)$ convergence
\[
 e^{-U/2}U^{1/2}H_U\mathbf1_U|_{(-X,X)}
 \longrightarrow
 -\sqrt2\,\operatorname{sgn}(u)I_0(|u|).
\tag{O3Z.2}\label{eq:o3z-boundary-functional}
\]
Hence the normalized functionals
\[
 \lambda_{X,U}(f):=e^{-U/2}U^{1/2}\ell_U(f)
\]
converge in $\mathcal H_X^*$ to
\[
 \lambda_{X,\infty}(f)=-\sqrt2\,\beta_X^{(0)}(f).
\tag{O3Z.3}\label{eq:o3z-functional-limit}
\]
Let $U_n\to\infty$ and $f_n\rightharpoonup f$ weakly in $\mathcal H_X$.
If $\beta_X^{(0)}(f)\ne0$, dual-norm convergence gives
$\lambda_{X,U_n}(f_n)\to-\sqrt2\beta_X^{(0)}(f)\ne0$.  Using the elementary
constant-mode lower bound and~\eqref{eq:mj-constant-norm},
\[
 a_{X,U_n}(f_n)
 \ge\frac{|\ell_{U_n}(f_n)|^2}{d_{U_n}}
 \gg_f\frac{e^{U_n}}{U_n^2}\longrightarrow\infty.
\tag{O3Z.4}\label{eq:o3z-outside-divergence}
\]
This is the required liminf outside $H_X^0$.

If $f\in H_X^0$, exact Gamma compatibility and positivity of the Schur term give
\[
 a_{X,U_n}(f_n)\ge\mathfrak c_{\Gamma,X}[f_n].
\]
At fixed $X$ the Gamma graph norm and the $X$-graph norm are equivalent by
\eqref{eq:graph-norm-equiv}; therefore the Gamma quadratic form is weakly lower
semicontinuous on $\mathcal H_X$, and
\[
 \liminf_n a_{X,U_n}(f_n)
 \ge\mathfrak c_{\Gamma,X}[f].
\]

For the recovery condition, first let
$f\in C_c^\infty((-X,X))_{\rm odd}\cap H_X^0$ be nonzero.  By
Corollary~\ref{cor:first-jet}, its first nonzero jet has some finite order $m>0$.
Choose a fixed smooth odd $f_0$ with first jet order zero and use the exact near-null
vector
\[
 z_U(f)
 :=f-\frac{\ell_U(f)}{\ell_U(f_0)}f_0.
\]
Proposition~\ref{prop:o3n-nearnull} gives $z_U(f)\to f$, while
Corollary~\ref{cor:o3p-gamma-limit} gives
\[
 a_{X,U}(z_U(f))\longrightarrow\mathfrak c_{\Gamma,X}[f].
\tag{O3Z.5}\label{eq:o3z-smooth-recovery}
\]
The zero vector is trivial.

For arbitrary $f\in H_X^0$, choose smooth odd $g_n\to f$ in the fixed graph norm by
Theorem~\ref{thm:smooth-odd-core}.  With the same $f_0$, put
\[
 h_n:=g_n-
 \frac{\beta_X^{(0)}(g_n)}{\beta_X^{(0)}(f_0)}f_0.
\]
Then $h_n$ is smooth odd, belongs to $H_X^0$, and tends to $f$.  Apply
\eqref{eq:o3z-smooth-recovery} to each $h_n$ and choose a standard diagonal sequence
$n=n(U)\to\infty$ sufficiently slowly.  This gives $f_U\to f$ with
\[
 a_{X,U}(f_U)\to\mathfrak c_{\Gamma,X}[f].
\]
Thus both Mosco conditions hold.
\end{proof}

Theorem~\ref{thm:o3z-mosco} does not contradict fixed-vector divergence: its recovery
vectors are terminal-dependent and use precisely the $O(U^{-m})$ correction exposed in
R17/R26.

Now fix $0<R<S<T_0$ and put, for $X\in\{R,S\}$,
\[
 B_X:=G_{X,T_0}^-,
 \qquad
 A_X(U):=B_X^{-1/2}G_{X,U}^-B_X^{-1/2}.
\]
Define the baseline-whitened effective hyperplane
\[
 V_X:=B_X^{1/2}H_X^0
\]
and the transported Gamma form
\[
 \widetilde{\mathfrak c}_{\Gamma,X}[x,y]
 :=\mathfrak c_{\Gamma,X}[B_X^{-1/2}x,B_X^{-1/2}y].
\]

\begin{corollary}[Whitened constrained Gamma limit]
\label{cor:o3z-whitened-mosco}
The relative metric forms
\[
 x\longmapsto\langle A_X(U)x,x\rangle
\]
Mosco-converge to
\[
 \boxed{
 \widetilde a_{X,\infty}(x)
 :=
 \begin{cases}
 \widetilde{\mathfrak c}_{\Gamma,X}[x],&x\in V_X,\\[2pt]
 +\infty,&x\notin V_X.
 \end{cases}
 }
\tag{O3Z.6}\label{eq:o3z-whitened-limit}
\]
Moreover
\[
 \widetilde{\mathfrak c}_{\Gamma,X}[x]
 \ge a_0\|x\|^2,
 \qquad
 a_0=(1+\|H_{T_0}\|^2)^{-1},
\tag{O3Z.7}\label{eq:o3z-limit-coercivity}
\]
with the same uniform constant as in Proposition~\ref{prop:o3w-coercivity}.
\end{corollary}

\begin{proof}
The fixed boundedly invertible map $B_X^{1/2}$ is a homeomorphism for both the strong
and weak Hilbert topologies, so Theorem~\ref{thm:o3z-mosco} is preserved under the
change of variables $x=B_X^{1/2}f$.  Finally
\eqref{eq:graph-norm-equiv} at the fixed terminal $T_0$ gives
\[
 \|x\|^2
 =q_{T_0}^X(J_{X,T_0}f)
 \le(1+\|H_{T_0}\|^2)\mathfrak c_{\Gamma,X}[f],
\]
which is~\eqref{eq:o3z-limit-coercivity}.
\end{proof}

\begin{theorem}[Strong constrained-resolvent limit]
\label{thm:o3z-resolvent-limit}
For every fixed $t\ge0$ there is a bounded positive operator
$\mathcal R_{X,\infty}(t)$ with range in $V_X$, characterized by
\[
 \boxed{
 \widetilde{\mathfrak c}_{\Gamma,X}
 [\mathcal R_{X,\infty}(t)x,v]
 +t\langle\mathcal R_{X,\infty}(t)x,v\rangle
 =\langle x,v\rangle
 \quad(v\in V_X).
 }
\tag{O3Z.8}\label{eq:o3z-limit-riesz}
\]
It satisfies
\[
 \|\mathcal R_{X,\infty}(t)\|\le(a_0+t)^{-1}
\]
and
\[
 \boxed{
 (A_X(U)+tI)^{-1}
 \xrightarrow[U\to\infty]{\rm strong}
 \mathcal R_{X,\infty}(t).
 }
\tag{O3Z.9}\label{eq:o3z-resolvent-strong}
\]
The statement includes $t=0$.
\end{theorem}

\begin{proof}
For fixed $x$, $(A_X(U)+tI)^{-1}x$ is the unique minimizer of
\[
 F_{U,x,t}(y)
 :=\langle A_X(U)y,y\rangle+t\|y\|^2
 -2\operatorname{Re}\langle x,y\rangle.
\]
Corollary~\ref{cor:o3z-whitened-mosco} identifies the Mosco limit with the corresponding
functional on $V_X$.  Coercivity~\eqref{eq:o3z-limit-coercivity} gives the unique
limiting minimizer characterized by~\eqref{eq:o3z-limit-riesz}.  The Mosco liminf and
a recovery sequence give convergence of minimum values.  Since the finite functionals
are uniformly $(a_0+t)$-strongly convex, comparison with a recovery sequence for the
limiting minimizer upgrades weak convergence of minimizers to strong convergence.  This
proves~\eqref{eq:o3z-resolvent-strong}.  The same coercivity makes the argument valid at
$t=0$.
\end{proof}

Because the effective domain $V_X$ is a proper closed hyperplane, the theorem is stated
as a constrained-Riesz limit and is not rebranded as the resolvent of a densely defined
selfadjoint operator on the whole standardized Hilbert space.

\begin{corollary}[Strong inverse-root limit and resolution of the R24/R25 existence gate]
\label{cor:o3z-inverse-root-limit}
Define
\[
 \boxed{
 T_{X,\infty}
 :=\frac1\pi\int_0^\infty
 t^{-1/2}\mathcal R_{X,\infty}(t)\,dt.
 }
\tag{O3Z.10}\label{eq:o3z-T-limit}
\]
Then
\[
 \boxed{
 A_X(U)^{-1/2}\xrightarrow[s]{}T_{X,\infty}.
 }
\tag{O3Z.11}\label{eq:o3z-inverse-root-strong}
\]
Consequently
\[
 \boxed{
 D_U^-:=A_S(U)^{-1/2}W-WA_R(U)^{-1/2}
 \xrightarrow[s]{}
 D_\infty^-:=T_{S,\infty}W-WT_{R,\infty},
 }
\tag{O3Z.12}\label{eq:o3z-D-limit}
\]
and the R25 positive compression defect satisfies
\[
 \boxed{
 \mathscr I_U=W^*D_U^-
 \xrightarrow[s]{}
 \mathscr I_\infty:=W^*D_\infty^-\ge0.
 }
\tag{O3Z.13}\label{eq:o3z-I-limit}
\]
\end{corollary}

\begin{proof}
By Proposition~\ref{prop:o3w-coercivity},
\[
 \|(A_X(U)+tI)^{-1}\|\le(a_0+t)^{-1}.
\]
The integrable bound
\[
 \frac{t^{-1/2}}{a_0+t}
\]
and Theorem~\ref{thm:o3z-resolvent-limit} allow vectorwise dominated convergence in
the inverse-square-root integral.  This proves
\eqref{eq:o3z-inverse-root-strong}.  The remaining assertions follow from the fixed
isometry $W$ and~\eqref{eq:o3x-I-compression}.  Positivity of the limit follows from
positivity of every $\mathscr I_U$.
\end{proof}

Thus the open part of R24/R25 is no longer the existence of an asymptotic inverse-root
orientation: the strong limit exists.  What remains is to decide whether the fixed
operator $D_\infty^-$ vanishes.

There is a fixed-window characterization on the effective higher-jet domain.  Put
\[
 V_0:=WV_R\subset V_S,
 \qquad
 K:=V_S\cap V_0^\perp,
\]
where orthogonality is in the baseline-whitened Hilbert metric.

\begin{proposition}[Fixed Gamma cross-block criterion]
\label{prop:o3z-gamma-block}
For $x\in V_R$ and one, equivalently every, fixed $t\ge0$,
\[
 \boxed{
 \left.
 \bigl(W^*\mathcal R_{S,\infty}(t)W
       -\mathcal R_{R,\infty}(t)\bigr)
 \right|_{V_R}=0
 \quad\Longleftrightarrow\quad
 \widetilde{\mathfrak c}_{\Gamma,S}[WV_R,K]=0.
 }
\tag{O3Z.14}\label{eq:o3z-gamma-block}
\]
If the cross block is nonzero, the limiting positive resolvent defect is nonzero on
$V_R$; if it vanishes, the constrained Gamma resolvents intertwine there.
\end{proposition}

\begin{proof}
Gamma compatibility and the definition of $W$ give
\[
 \widetilde{\mathfrak c}_{\Gamma,S}[Wx,Wy]
 =\widetilde{\mathfrak c}_{\Gamma,R}[x,y]
 \qquad(x,y\in V_R).
\tag{O3Z.15}\label{eq:o3z-gamma-compat}
\]
The limit Gamma form is bounded and coercive in the whitened baseline norm, hence is
represented on $V_S=V_0\oplus K$ by a bounded positive block operator
\[
 \begin{pmatrix}
 W L_RW^*&C^*\\
 C&D
 \end{pmatrix},
\]
where
\[
 \langle CWx,k\rangle
 =\widetilde{\mathfrak c}_{\Gamma,S}[Wx,k].
\]
For input in $V_0$, the upper-left block of the inverse after adding $tI$ is the inverse
of the Schur complement
\[
 W(L_R+tI)W^*-C^*(D+tI)^{-1}C.
\]
It equals $W(L_R+tI)^{-1}W^*$ identically iff
$C^*(D+tI)^{-1}C=0$, equivalently $C=0$.  This is precisely the stated Gamma
cross-orthogonality condition and is independent of $t$.
\end{proof}

\begin{remark}[R27 reduction and firewall]
\label{rem:o3z-firewall}
The $U\to\infty$ inverse-root existence problem has now been reduced to a fixed
$R,S,T_0$ Gamma block.  In original source variables the remaining question asks
whether
\[
 \mathfrak c_{\Gamma,S}[J_{R,S}f,h]
\]
vanishes for every $f\in\ker\beta_R^{(0)}$ and every
$h\in\ker\beta_S^{(0)}$ whose whitened representative lies in the baseline-orthogonal
complement of the nested $R$ hyperplane.

This is still a modulus/inverse-functional-calculus statement.  R14 shows that even
perfect modulus compatibility need not synchronize the polar gauges.  Therefore no
conclusion about $\Gamma_U$, the R22 angle defect, strong terminal transport, or a
global Object~X follows from the present Mosco limit alone.
\end{remark}

\begin{openproblem}[R27-F: fixed Gamma block]
\label{open:o3z-fixed-gamma-block}
Decide the fixed-window cross block in
\eqref{eq:o3z-gamma-block} for the concrete P11 baseline.  A nonzero witness would give
a persistent concrete inverse-root orientation defect.  Vanishing would prove
asymptotic inverse-root intertwining on the effective higher-jet domain.  Either outcome
would still require a separate polar-gauge analysis before any terminal-transport
conclusion.
\end{openproblem}
